\documentclass[11pt]{article}

% Shared preamble for Thurstone papers.
% Imported via % Shared preamble for Thurstone papers.
% Imported via % Shared preamble for Thurstone papers.
% Imported via \input{../shared/preamble.tex} from each paper.tex.
% Keep packages and macros common to all papers here.

\usepackage[utf8]{inputenc}
\usepackage[T1]{fontenc}
\usepackage{lmodern}
\usepackage{amsmath,amssymb,amsthm}
\usepackage{graphicx}
\graphicspath{{figures/}}
\usepackage{booktabs}
\usepackage{microtype}
\usepackage[hidelinks]{hyperref}

% Bibliography. We use biblatex with the bibtex backend because Tectonic runs
% bibtex natively (no external biber, which avoids biber/biblatex version
% mismatches). bibtex runs through Tectonic's own I/O, so the shared
% bibliography at papers/refs.bib resolves directly via the parent path below.
\usepackage[backend=bibtex,style=authoryear,sorting=nyt]{biblatex}
\addbibresource{../refs.bib}

% --- Shared macros ---
\newcommand{\thurstone}{Thurstone}
 from each paper.tex.
% Keep packages and macros common to all papers here.

\usepackage[utf8]{inputenc}
\usepackage[T1]{fontenc}
\usepackage{lmodern}
\usepackage{amsmath,amssymb,amsthm}
\usepackage{graphicx}
\graphicspath{{figures/}}
\usepackage{booktabs}
\usepackage{microtype}
\usepackage[hidelinks]{hyperref}

% Bibliography. We use biblatex with the bibtex backend because Tectonic runs
% bibtex natively (no external biber, which avoids biber/biblatex version
% mismatches). bibtex runs through Tectonic's own I/O, so the shared
% bibliography at papers/refs.bib resolves directly via the parent path below.
\usepackage[backend=bibtex,style=authoryear,sorting=nyt]{biblatex}
\addbibresource{../refs.bib}

% --- Shared macros ---
\newcommand{\thurstone}{Thurstone}
 from each paper.tex.
% Keep packages and macros common to all papers here.

\usepackage[utf8]{inputenc}
\usepackage[T1]{fontenc}
\usepackage{lmodern}
\usepackage{amsmath,amssymb,amsthm}
\usepackage{graphicx}
\graphicspath{{figures/}}
\usepackage{booktabs}
\usepackage{microtype}
\usepackage[hidelinks]{hyperref}

% Bibliography. We use biblatex with the bibtex backend because Tectonic runs
% bibtex natively (no external biber, which avoids biber/biblatex version
% mismatches). bibtex runs through Tectonic's own I/O, so the shared
% bibliography at papers/refs.bib resolves directly via the parent path below.
\usepackage[backend=bibtex,style=authoryear,sorting=nyt]{biblatex}
\addbibresource{../refs.bib}

% --- Shared macros ---
\newcommand{\thurstone}{Thurstone}


\newtheorem{theorem}{Theorem}
\newtheorem{proposition}[theorem]{Proposition}

\newcommand{\E}{\mathbb{E}}
\newcommand{\Prob}{\mathbb{P}}
\newcommand{\R}{\mathbb{R}}
\newcommand{\ind}{\mathbf{1}}

\title{All-Action Propensities for Correlated Gaussian Thompson Sampling}
\author{Peter Cotton}
\date{August 2026\\[0.5em]\normalsize\emph{Working manuscript --- empirical results pending}}

\begin{document}
\maketitle

\begin{abstract}
Thompson sampling is easy to execute: sample once from a posterior and choose the action with the largest sampled value. Its action probabilities can nevertheless be difficult to recover. For a finite action set, the logging propensity of action $i$ is the posterior winner probability
\[
p_i = \Prob\!\left( U_i = \max_j U_j \mid \mathcal{H}, x \right),
\]
where $U = (U_1, \ldots, U_N)$ is the jointly sampled utility vector. These propensities are required by inverse-propensity and doubly robust off-policy estimators. A familiar one-dimensional product formula is exact when the sampled action values are independent. It is generally invalid for a shared posterior parameter draw, because univariate Gaussian marginals do not determine winner probabilities.

We study correlated Gaussian Thompson sampling under the factor-plus-diagonal representation
\[
U_i = \mu_i + v_i^\top f + \sqrt{D_i}\,\epsilon_i,
\qquad
f \sim N(0, I_k),
\qquad
\epsilon_i \stackrel{\mathrm{iid}}{\sim} N(0,1),
\]
with $D_i > 0$. Conditional on the factors, the actions are independent. This yields a shared product-CDF field from which all $N$ winner probabilities can be evaluated together in $O(QN(L+k))$ arithmetic, where $Q$ is the number of factor nodes and $L$ is the size of a one-dimensional utility quadrature.

We prove that Gaussian marginals alone are insufficient, give the correct Gaussian orthant representation, establish continuity to the singular pure-factor boundary, and bound the error induced by a factor approximation through Gaussian total variation. We also describe the Jacobian of the propensity map as a weighted ``photo-finish'' graph Laplacian, providing all-action sensitivity to posterior-mean bonuses.

For off-policy evaluation, we derive exact conditional bias formulas when an approximate propensity vector is substituted into inverse-propensity or doubly robust estimators. The formulas show why small absolute errors on rare actions can create large counterfactual bias. The paper closes with a falsifiable computational and statistical experimental program. The intended claim is not that factor-Gaussian evaluation replaces direct posterior sampling in every Thompson-sampling model, but that it makes correlated all-action propensities explicit in an important structured regime where independence-based calculations are wrong and brute-force Monte Carlo can be costly.
\end{abstract}

\section{Introduction}

Thompson sampling \parencite{thompson1933,chapelle2011,russo2018} converts posterior uncertainty into randomized decisions by drawing one model or value function from the current posterior and selecting its maximizing action. In a contextual linear bandit \parencite{agrawal2013}, for example, a Gaussian parameter draw produces a vector of correlated action scores, and the maximizer is played.

Let $\mathcal{H}$ denote the information available immediately before a decision, let $x$ be the current context, and let $U = (U_1, \ldots, U_N)$ be the jointly sampled posterior values of the available actions. Assuming ties have probability zero, the Thompson-sampling logging propensity is
\[
p_i(x, \mathcal{H})
=
\Prob\!\left( i = \arg\max_{1 \le j \le N} U_j \mid x, \mathcal{H} \right).
\]
The vector $p = (p_1, \ldots, p_N)$ is both the posterior distribution of the optimizer and the conditional categorical policy generated by Thompson sampling.

The distinction between \emph{sampling an action} and \emph{evaluating all action probabilities} matters whenever logged Thompson-sampling data are used for counterfactual inference. Inverse-propensity scoring, doubly robust estimation, policy learning, exploration diagnostics, and support audits require the logging probability of the observed action, and sometimes the entire policy vector. General off-policy estimators are well developed once the logging policy is known \parencite{dudik2011,wang2017}. Under adaptive data collection, valid uncertainty quantification additionally requires procedures that respect the time dependence induced by the learning policy \parencite{hadad2021,zhan2021}.

Recent work has emphasized that Thompson-sampling propensities are not ordinarily exposed by the sampling operation and has derived useful expressions for several independent action-value families \parencite{jeunen2025}. For independent continuous action values with densities $f_i$ and CDFs $F_i$,
\begin{equation}
\label{eq:independent}
p_i = \int_{-\infty}^{\infty} f_i(u) \prod_{j \ne i} F_j(u)\, du.
\end{equation}
This identity is exact and computationally attractive in its proper domain.

The difficulty is that traditional parameter-sampling Thompson sampling uses a \emph{single shared posterior draw}. If $\widetilde\theta \sim N(m, S)$ and $U_i = z_i^\top \widetilde\theta$, then
\[
\operatorname{Cov}(U_i, U_j) = z_i^\top S z_j,
\]
which is generally nonzero. Retaining only the univariate distributions of the $U_i$ and inserting them into the independent-action expression discards precisely the dependence created by the common draw.

This paper studies the structured correlated case
\begin{equation}
\label{eq:factor}
U_i = \mu_i + v_i^\top f + \sqrt{D_i}\,\epsilon_i,
\qquad
f \sim N(0, I_k),
\qquad
\epsilon_i \stackrel{\mathrm{iid}}{\sim} N(0,1),
\qquad
D_i > 0.
\end{equation}
Consequently
\[
U \sim N(\mu, \Sigma),
\qquad
\Sigma = V V^\top + \operatorname{diag}(D).
\]
The representation is exact for Gaussian random-effects bandits with shared factors and action-specific posterior effects. It is also a natural approximation to finite Gaussian-process posteriors, Bayesian matrix-factorization recommenders, and correlated Bayesian last-layer models.

The computational primitive is a shared-field factor-probit evaluator developed in companion work \parencite{cotton2026probit}. Conditional on $f$, the action utilities are independent, but the same product of conditional CDFs serves every action. This converts $N$ separate winner calculations into one common field and $N$ readouts. The broader factor-probit program also supplies the potential, inverse map, Jacobian-vector products, and graph-Laplacian Hessian.

The present paper asks what this primitive contributes specifically to Thompson sampling and counterfactual inference. Its contributions are:

\begin{enumerate}
\item \textbf{Correlation correction.} We isolate the failure of marginal-only Gaussian propensity calculations under a shared posterior draw. In the binary case, one covariance term can change a propensity from $0.760$ to $0.987$ while leaving both univariate marginals unchanged.

\item \textbf{All-action factor-Gaussian evaluation.} We derive the conditional shared-field representation and an $O(QN(L+k))$ evaluator for the complete propensity vector. We distinguish deterministic low-dimensional quadrature from fixed-node, resampling-free quasi-Monte Carlo at moderate factor rank.

\item \textbf{Boundary and approximation guarantees.} We prove convergence to a pure-factor Thompson policy as the idiosyncratic variance tends to zero, provided the limiting posterior winner is almost surely unique. We also bound the propensity error caused by replacing a full Gaussian covariance by a factor-plus-diagonal approximation.

\item \textbf{Decision sensitivity.} We show that the Jacobian of $p$ with respect to posterior mean utilities is a weighted graph Laplacian. Its edges measure posterior density of pairwise ties at the winning boundary and identify the action pairs that actually compete for selection.

\item \textbf{Counterfactual consequences.} We derive the exact conditional bias of inverse-propensity and doubly robust estimators when the true correlated propensity $p$ is replaced by an approximation $\widetilde p$.
\end{enumerate}

The scope is deliberately limited. We do not claim a new definition of Thompson sampling, a new general theory of off-policy evaluation, or a universal algorithm for arbitrary Gaussian orthant probabilities. Nor do we claim that computing the full propensity vector is necessary merely to execute Thompson sampling. The target is narrower: many-action posterior utilities with low-rank shared dependence and non-negligible idiosyncratic uncertainty, especially when accurate logged propensities or their derivatives are required.

\section{Thompson propensities and the missing dependence}
\label{sec:missing}

\subsection{The posterior optimizer distribution}

At decision time $t$, write $\mathcal{F}_{t-1}$ for the history, $X_t$ for the current context, and $\mathcal{A}_t = \{1, \ldots, N_t\}$ for the available actions. A posterior-sampling policy draws a random value vector $U_t$ conditionally on $(\mathcal{F}_{t-1}, X_t)$ and selects
\[
A_t = \arg\max_{i \in \mathcal{A}_t} U_{t,i}.
\]
The conditional policy is therefore $p_{t,i} = \Prob(A_t = i \mid \mathcal{F}_{t-1}, X_t)$.

Suppressing the current conditioning, a general continuous joint density gives
\[
p_i
=
\int_{-\infty}^{\infty}
f_{U_i}(u)\,
\Prob\!\left( U_{-i} \le u \ind \mid U_i = u \right)
du.
\]
The independent expression \eqref{eq:independent} follows only when the conditional probability reduces to $\prod_{j \ne i} F_j(u)$. For a jointly Gaussian vector, zero cross-covariances are equivalent to independence. A shared Gaussian parameter draw usually produces neither.

\subsection{Exact Gaussian orthant representation}

For a fixed action $i$, let $B_i \in \R^{(N-1) \times N}$ have rows $e_j^\top - e_i^\top$ for $j \ne i$. Action $i$ wins exactly when $B_i U \le 0$ componentwise. If $U \sim N(\mu, \Sigma)$, then
\begin{equation}
\label{eq:orthant}
p_i
=
\Prob(B_i U \le 0)
=
\Phi_{N-1}\!\left( -B_i \mu;\; B_i \Sigma B_i^\top \right),
\qquad
\Phi_m(b; C) = \Prob(Z \le b),\quad Z \sim N(0, C).
\end{equation}
This formula is exact for arbitrary Gaussian dependence. It also makes the computational problem apparent: a straightforward all-action calculation calls an $(N-1)$-dimensional Gaussian probability routine $N$ times, with a different contrast covariance for every action. General-purpose methods for these probabilities include Genz transformations \parencite{genz1992} and Botev's minimax tilting \parencite{botev2017}.

\subsection{Univariate marginals do not identify propensities}

\begin{proposition}[Binary correlated-Gaussian propensity]
\label{prop:binary}
Let $(U_1, U_2)$ be jointly Gaussian with mean $(\mu_1, \mu_2)$ and covariance entries $\Sigma_{ab}$. Then
\[
p_1
=
\Prob(U_1 > U_2)
=
\Phi\!\left(
\frac{\mu_1 - \mu_2}
{\sqrt{\Sigma_{11} + \Sigma_{22} - 2\Sigma_{12}}}
\right).
\]
Consequently, the two univariate marginals do not determine $p_1$.
\end{proposition}

\begin{proof}
The difference $U_1 - U_2$ is Gaussian with mean $\mu_1 - \mu_2$ and variance $\Sigma_{11} + \Sigma_{22} - 2\Sigma_{12}$. The probability follows immediately. Holding $\mu_1, \mu_2, \Sigma_{11}$, and $\Sigma_{22}$ fixed while varying the admissible covariance $\Sigma_{12}$ leaves both marginals unchanged but changes the distribution of the difference.
\end{proof}

\subsection{Numerical example}
\label{sec:numexample}

Take $\mu_1 - \mu_2 = 1$, $\Sigma_{11} = \Sigma_{22} = 1$, and $\rho = \Sigma_{12}$. The independent-marginal calculation always returns $\Phi(1/\sqrt{2}) \approx 0.76025$. The correct probabilities are:

\begin{center}
\begin{tabular}{rrr}
\toprule
Correlation $\rho$ & Correct $p_1$ & Independent-marginal value \\
\midrule
$-0.9$ & $0.69602$ & $0.76025$ \\
$-0.5$ & $0.71815$ & $0.76025$ \\
$0$    & $0.76025$ & $0.76025$ \\
$0.5$  & $0.84134$ & $0.76025$ \\
$0.9$  & $0.98733$ & $0.76025$ \\
\bottomrule
\end{tabular}
\end{center}

At $\rho = 0.9$, the correct propensity is approximately $0.9873$, while the independence calculation is approximately $0.7603$. The discrepancy is structural, not a numerical-integration error.

\subsection{Relation to existing independent-action work}

\textcite{jeunen2025} derives the independent product formula \eqref{eq:independent} and its Gaussian specialization for independent action-value distributions. The same paper subsequently maps a shared linear Gaussian parameter posterior to the univariate marginal distribution of each action and states that these distributions can be inserted into the independent Gaussian formula.

Proposition~\ref{prop:binary} shows that this extension is not generally valid for a single shared parameter draw. The correct winner probability depends on the cross-action covariance.

This correction should not be overstated. The released implementation used for the empirical Open Bandit Pipeline experiment maintains a separate logistic model for each action, which is compatible with the independent-action setting. The issue concerns the broader shared-parameter interpretation, not the independent-action derivation or that particular empirical construction.

\section{Factor-Gaussian Thompson sampling}
\label{sec:factor}

\subsection{A canonical hierarchical model}

A simple random-effects posterior produces the factor-plus-diagonal model exactly. Let $\theta \sim N(m, S)$ and let $\alpha_i \sim N(a_i, D_i)$ independently across $i$ and independently of $\theta$. Suppose the sampled value of action $i$ is $U_i = z_i^\top \theta + \alpha_i$. Factor $S = L L^\top$ and write $\theta = m + L f$ with $f \sim N(0, I_k)$. Then
\[
U_i
=
z_i^\top m + a_i
+
(L^\top z_i)^\top f
+
\sqrt{D_i}\,\epsilon_i,
\]
so that
\[
\mu_i = z_i^\top m + a_i,
\qquad
v_i = L^\top z_i,
\qquad
\Sigma = V V^\top + \operatorname{diag}(D).
\]
The shared factors encode posterior uncertainty common to many actions; the diagonal terms encode action-specific posterior uncertainty. The same algebra applies whenever a finite Gaussian value posterior is represented or approximated by $\Sigma \approx V V^\top + \operatorname{diag}(D)$.

\subsection{The shared conditional field}

Conditionally on a factor realization $f$, define $m_i(f) = \mu_i + v_i^\top f$. Then $U_i \mid f \sim N(m_i(f), D_i)$ independently across $i$. Let $F_i(x \mid f)$ and $g_i(x \mid f)$ denote the corresponding conditional Gaussian CDF and density.

\begin{theorem}[Shared-field propensity representation]
\label{thm:sharedfield}
Under the factor-Gaussian model \eqref{eq:factor},
\[
p_i
=
\E_f \int_{-\infty}^{\infty}
g_i(x \mid f)
\prod_{j \ne i} F_j(x \mid f)\, dx.
\]
Define the common field
\[
H_f(x) = \prod_{j=1}^N F_j(x \mid f).
\]
Then every action probability is obtained from that field:
\[
p_i
=
\E_f \int_{-\infty}^{\infty}
g_i(x \mid f)\,
\frac{H_f(x)}{F_i(x \mid f)}\, dx.
\]
The quotient is evaluated numerically in the log domain, with the original product defining the limiting integrand in extreme tails.
\end{theorem}

\begin{proof}
Condition on $f$ and on $U_i = x$. Conditional independence gives
\[
\Prob(U_j \le x\ \forall j \ne i \mid U_i = x, f)
=
\prod_{j \ne i} F_j(x \mid f).
\]
Integrating against $g_i(x \mid f)$ and averaging over $f$ gives the first formula. Multiplying and dividing by $F_i(x \mid f)$ gives the shared-field form.
\end{proof}

The important computational fact is not merely that conditioning reduces the calculation to a one-dimensional integral. It is that the expensive product is shared by all $N$ alternatives.

\section{All-action evaluation}
\label{sec:allaction}

\subsection{Fixed-node discretization}

Let $\{(f_q, a_q)\}_{q=1}^Q$ be factor nodes and weights approximating expectation under $N(0, I_k)$, and let $\{(x_\ell, b_\ell)\}_{\ell=1}^L$ approximate the scalar integral.

Product Gauss--Hermite nodes are deterministic and highly accurate when $k$ is very small. For moderate $k$, fixed Sobol or randomized quasi-Monte Carlo designs are more practical. In that regime the proper description is \emph{fixed-node} or \emph{resampling-free}, rather than exact deterministic quadrature.

For each factor node and scalar point, compute
\[
z_{qi\ell} = \frac{x_\ell - \mu_i - v_i^\top f_q}{\sqrt{D_i}},
\qquad
\log F_{qi\ell} = \log \Phi(z_{qi\ell}),
\qquad
\log g_{qi\ell} = -\tfrac12 z_{qi\ell}^2 - \tfrac12 \log(2\pi D_i),
\]
and
\[
\log H_{q\ell} = \sum_{j=1}^N \log F_{qj\ell}.
\]
The numerical propensity is
\[
\widehat p_i
=
\sum_{q=1}^Q a_q
\sum_{\ell=1}^L b_\ell
\exp\!\left(
\log g_{qi\ell} + \log H_{q\ell} - \log F_{qi\ell}
\right).
\]
A stable implementation uses a high-quality $\log \Phi$ function, blocks the scalar lattice, and never forms $H / F_i$ in ordinary floating-point arithmetic.

\subsection{Algorithm}

\begin{enumerate}
\item Choose factor nodes $(f_q, a_q)$ and scalar integration nodes $(x_\ell, b_\ell)$.
\item For each factor node $q$, compute all conditional means $m_i(f_q)$.
\item For each scalar block, compute all conditional log CDFs and log densities.
\item Reduce over actions once to obtain the shared log product $\log H_{q\ell} = \sum_j \log F_{qj\ell}$.
\item Read out every action integrand by subtracting its own $\log F_{qi\ell}$ and adding its log density.
\item Integrate over the scalar lattice and average over factor nodes.
\item Report the probability-mass defect $\bigl| 1 - \sum_i \widehat p_i \bigr|$.
\end{enumerate}

If a final normalized probability vector is returned, the pre-normalization mass defect should remain a numerical diagnostic.

\begin{proposition}[Complexity]
\label{prop:complexity}
The complete propensity vector is evaluated in $O(QN(k+L))$ arithmetic. When $k \ll L$, this is $O(QNL)$. With scalar blocks of width $L_b$, working memory can be limited to $O(N L_b + N)$.
\end{proposition}

This is an all-action complexity result, not a universal superiority result. The factor-node budget may grow rapidly with factor rank, and direct Thompson sampling still obtains a single action with one posterior draw and one argmax.

\subsection{Comparison with direct Monte Carlo}

Direct posterior simulation draws $M$ utility vectors and estimates
\[
\widehat p_i^{\mathrm{MC}}
=
\frac{1}{M} \sum_{m=1}^M
\ind\!\left\{ i = \arg\max_j U_j^{(m)} \right\}.
\]
This estimator is unbiased and easily parallelized. Its standard deviation is $\sqrt{p_i(1 - p_i)/M}$. Reliable relative estimation of a rare probability requires $M p_i$ to be large. This is exactly the regime in which off-policy importance weights are sensitive.

The shared-field evaluator trades simulation variance for factor-integration and scalar-discretization error. The meaningful comparisons are therefore equal wall time, equal relative error, and equal downstream OPE error. Equal numbers of samples and quadrature nodes are not intrinsically meaningful.

\subsection{The pure-factor boundary}

Traditional shared-parameter linear Thompson sampling may produce $D = 0$. The smooth shared-field formula assumes $D_i > 0$, but it has a limiting relation to the singular model.

\begin{proposition}[Continuity to zero idiosyncratic variance]
\label{prop:continuity}
Let $U^{(\eta)} = \mu + V f + \eta \epsilon$ for $\eta > 0$. Suppose the maximizer of $U^{(0)} = \mu + V f$ is unique almost surely. Then
\[
\lim_{\eta \downarrow 0}
\Prob\!\left( i = \arg\max_j U_j^{(\eta)} \right)
=
\Prob\!\left( i = \arg\max_j U_j^{(0)} \right).
\]
\end{proposition}

\begin{proof}
Couple all $U^{(\eta)}$ using the same $f$ and $\epsilon$. Then $U^{(\eta)} \to U^{(0)}$ almost surely. At every realization for which $U^{(0)}$ has a unique maximizer, the argmax indicator is locally constant around $U^{(0)}$. The indicators therefore converge almost surely, and dominated convergence gives the result.
\end{proof}

If exact ties have positive probability, the limit depends on the vanishing-noise tie-breaking convention. Numerically, the $\eta \downarrow 0$ regime is also difficult because the conditional Gaussian CDF transitions become sharp.

\subsection{Error from covariance approximation}

Suppose the intended posterior is $P = N(\mu, \Sigma)$ and the factor approximation is $\widetilde P = N(\mu, \widetilde\Sigma)$. Let $p$ and $\widetilde p$ be their respective argmax distributions.

\begin{proposition}[Data-processing bound]
\label{prop:dataprocessing}
\[
\tfrac12 \|p - \widetilde p\|_1
\le
d_{\mathrm{TV}}(P, \widetilde P)
\le
\sqrt{ \tfrac12 D_{\mathrm{KL}}(P \| \widetilde P) }.
\]
For equal means,
\[
D_{\mathrm{KL}}(P \| \widetilde P)
=
\tfrac12
\left[
\operatorname{tr}(\widetilde\Sigma^{-1} \Sigma) - N
+ \log \frac{\det \widetilde\Sigma}{\det \Sigma}
\right].
\]
\end{proposition}

\begin{proof}
The argmax is a measurable map from $\R^N$ to the finite action set, so total variation contracts under the map. Pinsker's inequality gives the second bound.
\end{proof}

The bound separates posterior-factorization error from numerical integration error. It also makes clear that accurate approximation of marginal variances is insufficient: the choice-relevant contrast covariance must be approximated.

\section{Sensitivity and the photo-finish graph}
\label{sec:photofinish}

Define the Gaussian perturbed-maximum potential
\[
G(\mu) = \E \max_i U_i,
\qquad
U \sim N(\mu, \Sigma),
\]
with $\Sigma$ fixed. Since ties occur with probability zero when the diagonal variance is positive, $\nabla_\mu G(\mu) = p(\mu)$. Thus the propensity Jacobian is $J = \nabla_\mu^2 G$.

For $i \ne j$, define
\[
w_{ij}
=
\E_f \int_{-\infty}^{\infty}
g_i(x \mid f)\, g_j(x \mid f)
\prod_{\ell \ne i,j} F_\ell(x \mid f)\, dx.
\]
This is the posterior density of an $i$/$j$ tie at the winning boundary, integrated over the common factors.

\begin{proposition}[Photo-finish Laplacian]
\label{prop:laplacian}
\[
J = \sum_{i < j} w_{ij}\, (e_i - e_j)(e_i - e_j)^\top.
\]
Equivalently, $J_{ij} = -w_{ij}$ for $i \ne j$ and $J_{ii} = \sum_{j \ne i} w_{ij}$. Consequently $J \succeq 0$, $J \ind = 0$, and a small posterior-mean adjustment satisfies $dp = J\, d\mu$.
\end{proposition}

\begin{proof}
For $j \ne i$, differentiate the conditional propensity formula with respect to $\mu_j$. Since $\partial F_j(x \mid f) / \partial \mu_j = -g_j(x \mid f)$, we obtain $\partial p_i / \partial \mu_j = -w_{ij}$. A common shift of every utility leaves the winner unchanged, so every row of $J$ sums to zero. This determines the diagonal entries and gives the Laplacian representation.
\end{proof}

The graph with conductances $w_{ij}$ is a posterior competition graph. A large edge does not merely mean that two posterior values are correlated. It means they frequently meet near a tie while all other actions lie below them.

For deterministic posterior-mean bonuses $b$,
\[
p(\mu + b) = p(\mu) + J b + o(\|b\|).
\]
The graph therefore identifies which bonuses can meaningfully redirect exploration.

For fixed positive idiosyncratic variance, the companion factor-probit theory \parencite{cotton2026probit} establishes that the centered map
\[
\mu \in \ind^\perp \longmapsto p(\mu) \in \operatorname{int} \Delta_N
\]
is a smooth bijection. Hence a target interior action distribution $q$ can be implemented by solving for a centered utility adjustment producing $p = q$. This is an inverse probability-matching control layer. It deliberately changes the policy and should not be confused with unmodified posterior sampling.

\section{Counterfactual inference with correlated Thompson logs}
\label{sec:counterfactual}

\subsection{Contextual-bandit estimand}

At time $t$, let $\mathcal{F}_{t-1}$ be the observed past, $X_t$ the current context, $A_t$ the logged action, and $Y_t(a)$ the potential reward of action $a$. Let
\[
p_{t,a} = \Prob(A_t = a \mid \mathcal{F}_{t-1}, X_t)
\]
be the true correlated Thompson propensity, and let $\pi_{t,a}$ be a target policy measurable with respect to the same pre-action information. Define $m_{t,a} = \E[Y_t(a) \mid \mathcal{F}_{t-1}, X_t]$.

Assume sequential randomization,
\[
A_t \perp \{Y_t(a) : a \in \mathcal{A}_t\} \mid (\mathcal{F}_{t-1}, X_t),
\]
and overlap: $p_{t,a} > 0$ whenever $\pi_{t,a} > 0$. The one-step target value is
\[
v_t^\pi = \sum_a \pi_{t,a}\, m_{t,a}.
\]
This is a contextual-bandit estimand on the observed history and context distribution. It does not automatically evaluate the entire counterfactual trajectory of a different adaptive learning algorithm whose earlier actions would have changed its later posterior.

\subsection{Exact-propensity IPS and DR}

The inverse-propensity contribution is
\[
\widehat v_t^{\mathrm{IPS}} = \frac{\pi_{t,A_t}}{p_{t,A_t}}\, Y_t.
\]
For any predictable reward estimate $\widehat m_{t,a}$, define
\[
\widehat v_t^{\mathrm{DR}}
=
\sum_a \pi_{t,a}\, \widehat m_{t,a}
+
\frac{\pi_{t,A_t}}{p_{t,A_t}}
\left( Y_t - \widehat m_{t,A_t} \right).
\]

\begin{proposition}[Conditional unbiasedness]
\label{prop:unbiased}
Under sequential randomization and overlap,
\[
\E[\widehat v_t^{\mathrm{IPS}} \mid \mathcal{F}_{t-1}, X_t] = v_t^\pi
\qquad\text{and}\qquad
\E[\widehat v_t^{\mathrm{DR}} \mid \mathcal{F}_{t-1}, X_t] = v_t^\pi.
\]
\end{proposition}

\begin{proof}
For IPS,
\[
\E\!\left[ \frac{\pi_{t,A_t}}{p_{t,A_t}} Y_t \right]
=
\sum_a p_{t,a}\, \frac{\pi_{t,a}}{p_{t,a}}\, m_{t,a}
=
\sum_a \pi_{t,a}\, m_{t,a}.
\]
For DR, the conditional expectation of the residual correction is $\sum_a \pi_{t,a} (m_{t,a} - \widehat m_{t,a})$, which cancels the corresponding difference in the direct term.
\end{proof}

Adaptive collection does not destroy this conditional unbiasedness, but it does invalidate na\"ive i.i.d.\ uncertainty calculations. Martingale confidence intervals, stabilized estimators, or adaptive weighting should be used \parencite{hadad2021,zhan2021}.

\subsection{Bias from approximate propensities}

Suppose the evaluator supplies $\widetilde p_{t,a} > 0$ instead of the true $p_{t,a}$.

\begin{proposition}[Exact conditional bias under propensity error]
\label{prop:bias}
If IPS uses $\widetilde p$, then
\[
\E[\widetilde v_t^{\mathrm{IPS}}] - v_t^\pi
=
\sum_a \pi_{t,a}\, m_{t,a}
\left( \frac{p_{t,a}}{\widetilde p_{t,a}} - 1 \right).
\]
If DR uses $\widetilde p$, then
\[
\E[\widetilde v_t^{\mathrm{DR}}] - v_t^\pi
=
\sum_a \pi_{t,a}
\left( m_{t,a} - \widehat m_{t,a} \right)
\left( \frac{p_{t,a}}{\widetilde p_{t,a}} - 1 \right).
\]
If $|m_{t,a}| \le R_{\max}$, then
\[
\left| \E[\widetilde v_t^{\mathrm{IPS}}] - v_t^\pi \right|
\le
R_{\max} \sum_a \pi_{t,a}\,
\frac{|p_{t,a} - \widetilde p_{t,a}|}{\widetilde p_{t,a}}.
\]
\end{proposition}

\begin{proof}
For IPS,
\[
\E[\widetilde v_t^{\mathrm{IPS}}]
=
\sum_a p_{t,a}\, \frac{\pi_{t,a}}{\widetilde p_{t,a}}\, m_{t,a}.
\]
Subtracting the target value gives the first identity. For DR, take the expectation of the direct term and residual correction:
\[
\sum_a \pi_{t,a}\, \widehat m_{t,a}
+
\sum_a p_{t,a}\, \frac{\pi_{t,a}}{\widetilde p_{t,a}} (m_{t,a} - \widehat m_{t,a}).
\]
Subtracting $\sum_a \pi_{t,a} m_{t,a}$ gives the second identity. The bound follows by the triangle inequality.
\end{proof}

This result explains why absolute probability error is not an adequate metric for propensity evaluation. An error of $10^{-4}$ is negligible when the true propensity is $0.2$, but potentially severe when the propensity is $2 \times 10^{-4}$. Propensity methods intended for OPE should therefore report relative error, log-propensity error, and error stratified by true probability.

\subsection{OPE distortion in the binary example}

Return to the binary example of Section~\ref{sec:numexample} with $\rho = 0.9$. The true logging probability of action 1 is $p_1 \approx 0.987326$, while the independence approximation is $\widetilde p_1 \approx 0.760250$. Suppose the target policy always selects action 1. An IPS estimator using the incorrect propensity has conditional expectation
\[
\frac{p_1}{\widetilde p_1}\, m_1 \approx 1.29869\, m_1.
\]
The omitted covariance therefore produces approximately $29.9\%$ upward bias in this simple counterfactual calculation.

\subsection{Logging and reconstruction}

The safest operational design is to compute and log: the realized action's propensity; the posterior or model version; the available action set; the tie-breaking convention; and numerical mass-defect and convergence diagnostics.

Retrospective reconstruction requires the exact posterior state, feature transformation, action availability, and numerical convention that were in force. A fixed-node evaluator removes fresh simulation noise from repeated reconstructions, but it does not remove posterior misspecification.

\section{Experimental program}
\label{sec:experiments}

The analytical claims can be checked immediately. Practical value depends on an adversarial empirical program that separates numerical evaluation, statistical consequences, and end-to-end system cost.

\subsection{Operator accuracy and scaling}

Construct synthetic Gaussian races over $N \in \{10, 100, 1000, 10000\}$ and $k \in \{0, 1, 2, 4, 8\}$, with controlled mean spread, common-factor strength, diagonal variance heterogeneity, near-duplicate actions, and rare-action probabilities.

Compare: (i) shared-field product Gauss--Hermite quadrature for small $k$; (ii) fixed Sobol and randomized quasi-Monte Carlo factor nodes; (iii) direct hard-argmax Monte Carlo; (iv) Genz-type orthant calculations; (v) minimax-tilting estimates where feasible; and (vi) the independent-marginal approximation. Use Monte Carlo budgets $M \in \{10^2, 10^3, 10^4, 10^5, 10^6\}$.

Report maximum absolute error, total variation, relative error by probability bucket, log-propensity error, mass defect, wall time, peak memory, and run-to-run variability. Comparisons must be made at equal wall time and equal error, not merely equal node or sample counts.

\subsection{Correlation stress test}

Hold every univariate marginal fixed while varying only the shared correlation structure. Plot $p_i^{\mathrm{true}}$, $p_i^{\mathrm{factor}}$, and $p_i^{\mathrm{independent}}$ against correlation strength.

Extend the binary experiment to blocks of exchangeable actions; mixed positive and negative loadings; duplicated or near-duplicated actions; and strongly correlated actions that are far below the winning boundary. The photo-finish weights should predict which covariance changes materially affect the propensity vector.

\subsection{Off-policy bias and coverage}

Simulate a contextual random-effects bandit:
\[
\theta_t \mid \mathcal{F}_{t-1} \sim N(m_t, S_t),
\qquad
\alpha_{t,i} \mid \mathcal{F}_{t-1} \sim N(a_{t,i}, D_{t,i}),
\qquad
U_{t,i} = z_{t,i}^\top \theta_t + \alpha_{t,i}.
\]
Select the Thompson winner and generate rewards from a known outcome model.

Vary the number of actions, posterior factor rank, factor strength, target-policy overlap, reward-model accuracy, prevalence of rare actions, and posterior update frequency. Evaluate IPS, self-normalized IPS, DR, and adaptive-weighting estimators using: (i) high-accuracy true propensities; (ii) the shared-field approximation; (iii) finite-sample argmax Monte Carlo; and (iv) independent marginals. Report bias, RMSE, interval coverage, effective sample size, and empirical versus theoretically implied bias.

\subsection{Reproduce the independent case, then add sharing}

First reproduce an independent per-action Gaussian or logistic Thompson environment. In this regime, the shared-field evaluator with $V = 0$ must agree with the independent product formula.

Then introduce a shared Gaussian final-layer parameter while preserving the marginal means and variances. This creates a nested comparison: independent action posteriors within a shared correlated posterior. The independent formula should remain correct in the first regime and fail only when shared covariance is introduced.

\subsection{The $D \downarrow 0$ boundary}

For pure linear or Gaussian-process posterior draws, compare: (i) direct factor-space argmax Monte Carlo at $D = 0$; (ii) shared-field continuation with $D_i = \eta^2$; (iii) polyhedral Gaussian integration in factor space for small $k$; and (iv) action-difference orthant calculations for small $N$. Measure convergence, conditioning, and runtime as $\eta \downarrow 0$.

\subsection{Semi-synthetic recommendation}

Use a real context-action feature corpus but generate rewards from a known hierarchical Bayesian model. This preserves realistic action embeddings and availability patterns while retaining a known counterfactual value.

A natural construction is item recommendation with shared user uncertainty, item embeddings, item-specific random effects, and a finite changing candidate set. The principal outcome is not online regret. It is whether correlated propensities improve offline policy ranking, value estimation, and confidence-interval coverage at a practical logging cost.

\subsection{Success and pivot criteria}

The project succeeds if at least one of the following holds: rare-action relative error is materially lower at fixed wall time; OPE bias or interval coverage materially improves; independence causes operationally meaningful errors under realistic shared posteriors; the photo-finish graph predicts sensitivity or enables targeted approximations; or inverse probability matching provides useful policy control.

The project should pivot if: low factor rank is too restrictive; direct common-random-number Monte Carlo dominates at every relevant accuracy; near-singular $D$ covers most important models and is numerically prohibitive; or propensity error is negligible relative to ordinary overlap and reward-model error.

\section{Discussion}
\label{sec:discussion}

\subsection{When all-action evaluation is worthwhile}

Computing $p$ is not needed to take a Thompson action. One posterior sample and one argmax are cheaper.

The evaluator is justified when: propensities must be logged for OPE; rare-action probabilities require more accuracy than feasible argmax counts; the complete optimizer distribution or entropy is itself a diagnostic; derivatives with respect to posterior means or covariance are required; target exploration shares are imposed through inverse probability matching; or posterior snapshots permit batched offline reconstruction.

Direct Monte Carlo is likely preferable when only the selected action is required, factor rank is high, or rough propensities are adequate.

\subsection{Correlation is not merely a calibration detail}

A shared posterior parameter does more than alter a standard error. It changes the events defining action selection.

Highly positively correlated actions move together, so their difference can be much less uncertain than either marginal value. Negative correlation can make their difference more uncertain. Winner probabilities therefore live in \emph{contrast space}, not in the list of marginal means and variances.

This also clarifies action duplication. Two nearly identical actions with strongly correlated posterior values should not receive the same multiplicity benefit as two independent draws with identical marginals. An independence approximation can therefore overstate the policy effect of cloned or redundant candidates.

\subsection{Posterior, factorization, and numerical error}

Three errors must be separated: \emph{posterior-model error}, where the Gaussian or hierarchical posterior is misspecified; \emph{factorization error}, where $V V^\top + \operatorname{diag}(D)$ imperfectly approximates the intended covariance; and \emph{numerical error}, where finite factor nodes and scalar quadrature imperfectly evaluate the factor model. The total-variation proposition addresses the second in principle. Numerical benchmarks address the third. Neither resolves the first.

\subsection{Limitations}

Low factor rank is essential. Tensor quadrature is practical only for very small $k$, and fixed-node quasi-Monte Carlo may still require a large node count as rank increases.

The smooth shared-field algorithm assumes positive diagonal variance. Pure shared-parameter Thompson sampling lies on a singular boundary and may require a separate polyhedral method.

Exact propensities do not solve weak overlap. Correct importance weights can still have enormous variance.

The contextual-bandit estimand does not automatically evaluate a fully counterfactual adaptive trajectory.

Finally, the all-action algorithm has not yet been shown to dominate modern Monte Carlo or Gaussian-probability software in wall-clock performance. That is an empirical question, not a theorem.

\section{Conclusion}

Thompson sampling hides a categorical policy inside a posterior argmax. When action values are independent, a product of marginal CDFs exposes that policy. When they arise from a shared Gaussian posterior draw, the policy depends on cross-action covariance and cannot be reconstructed from univariate marginals.

For factor-plus-diagonal Gaussian value posteriors, conditional independence restores tractability without discarding dependence. A single shared product-CDF field yields all action propensities, their posterior-mean sensitivities, and a natural competition graph.

These quantities feed directly into counterfactual estimators, where propensity-ratio errors have exact and potentially large bias consequences.

The claim is therefore precise: structured Gaussian dependence turns Thompson propensities into a factor-probit problem, and the shared-field construction makes that problem computationally explicit. Whether this method is practically superior to large-sample simulation is a falsifiable question. The proposed experiments are designed to answer it.

\printbibliography

\end{document}
